\documentclass[11pt]{article}
\usepackage[a4paper,margin=1in]{geometry}
\usepackage{amsmath,amssymb}
\usepackage{booktabs,longtable}
\usepackage{graphicx}
\usepackage[hyphens]{url}
\usepackage[colorlinks=true,linkcolor=black,urlcolor=blue,citecolor=blue]{hyperref}
\setlength{\parskip}{0.55em}\setlength{\parindent}{0pt}
\providecommand{\tightlist}{\setlength{\itemsep}{0.2em}\setlength{\parskip}{0pt}}

\title{Supplementary Information — Forecast skill of structural models for Northern cod}
\author{\textbf{Amin Abaee}\\ Independent Researcher\\[2pt] ORCID: \href{https://orcid.org/0000-0002-0019-1842}{0000-0002-0019-1842}\\ Email: \href{mailto:amin_abaee@ut.ac.ir}{amin_abaee@ut.ac.ir}}
\date{September 18, 2026}
\begin{document}
\maketitle

\section{Supplementary Information
(v2)}\label{supplementary-information-v2}

\subsection{Forecast skill of structural models for Northern cod: a
scored
ladder}\label{forecast-skill-of-structural-models-for-northern-cod-a-scored-ladder}

This supplement supports the main article. Section numbers are cited in
the main text as SI-1, SI-2, and so on. All content here is derived from
the analysis code and result files archived with this study; no result
appears in this supplement that is not reproducible from that code.

\begin{center}\rule{0.5\linewidth}{0.5pt}\end{center}

\subsection{SI-1 Order in which the passes were
run}\label{si-1-order-in-which-the-passes-were-run}

The main text (Section 4) states that the study is a fixed computational
protocol rather than a prospective registration, that no dated
pre-scoring protocol file exists, and that the additional passes were
declared in the text as they were added. This section records the order
of execution, documenting which choices were fixed before scoring and
which were not.

The order below is established from the analysis code and its recorded
metadata, not from recollection. Passes 1 and 2 are labelled as such in
the driver script \texttt{run\_ladder.py} and are distinguished in the
archived \texttt{results/meta.json} by the keys \texttt{catch\_pass1}
and \texttt{catch\_pass2}. Passes 3 and 4 are separate scripts that
import from the pass-1/2 driver \texttt{run\_ladder.py}, which fixes
their order after it: \texttt{run\_xte.py} imports the model definitions
(\texttt{SPECS}) themselves, while the capelin scripts import the shared
machinery (\texttt{DATA}, \texttt{OUT}, \texttt{EPS}, \texttt{step},
\texttt{naive\_baselines}) and add one further model of their own,
\texttt{M\_cap\_index}.

{\def\LTcaptype{none} % do not increment counter
\begin{longtable}[]{@{}
  >{\raggedright\arraybackslash}p{(\linewidth - 8\tabcolsep) * \real{0.2000}}
  >{\raggedright\arraybackslash}p{(\linewidth - 8\tabcolsep) * \real{0.2000}}
  >{\raggedright\arraybackslash}p{(\linewidth - 8\tabcolsep) * \real{0.2000}}
  >{\raggedright\arraybackslash}p{(\linewidth - 8\tabcolsep) * \real{0.2000}}
  >{\raggedright\arraybackslash}p{(\linewidth - 8\tabcolsep) * \real{0.2000}}@{}}
\toprule\noalign{}
\begin{minipage}[b]{\linewidth}\raggedright
\#
\end{minipage} & \begin{minipage}[b]{\linewidth}\raggedright
Pass
\end{minipage} & \begin{minipage}[b]{\linewidth}\raggedright
Catch / data treatment
\end{minipage} & \begin{minipage}[b]{\linewidth}\raggedright
Produced by
\end{minipage} & \begin{minipage}[b]{\linewidth}\raggedright
Archived output
\end{minipage} \\
\midrule\noalign{}
\endhead
\bottomrule\noalign{}
\endlastfoot
1 & Regime catch & Regime constants 240/120/5 kt, read from the SAR
prose & \texttt{run\_ladder.py} (\texttt{\#\ Pass\ 1\ (regime\ C)}) &
\texttt{fixed\_window\_scores.csv}, \texttt{rolling\_summary.csv} (rows
\texttt{catch\ =\ regime}) \\
2 & Annual landings & Annual series, Schijns et al.~(2021) Table 1,
tonnes/1000 & \texttt{run\_ladder.py}
(\texttt{\#\ Pass\ 2:\ annual\ catch}) & same files, rows
\texttt{catch\ =\ annual} \\
2b & Survey start & Forecast initiated at the survey state, annual catch
& \texttt{run\_survey\_start()}, called within the pass-2 block of
\texttt{run\_ladder.py} & \texttt{survey\_start\_forecasts.csv} \\
3 & Capelin & Continuous capelin-W ablation, observed acoustic years
only & \texttt{run\_capelin\_index.py}, \texttt{run\_capelin\_regime.py}
& \texttt{capelin\_index\_summary.csv},
\texttt{capelin\_regime\_summary.csv} \\
4 & Specification B & The ladder on the xteNCAM domain alone,
\textbf{not pooled} with NCAM 2016 & \texttt{run\_xte.py} &
\texttt{xte\_rolling\_summary.csv},
\texttt{xte\_fixed\_window\_scores.csv} \\
\end{longtable}
}

Three points follow from this order and bear on how the scores should be
read.

\textbf{The five ladder models were fixed before any pass was scored.}
The \texttt{SPECS} list in \texttt{run\_ladder.py} defines M1, M1b, M2,
M3 and M4 once, and Specification B (pass 4) imports that same list
rather than redefining it. None of M1--M4 was altered after a score was
seen. The capelin pass is the one place where a model was added after
scoring began: it introduces \texttt{M\_cap\_index}, a prey-informed
variant, which is reported in the main text as its own ablation and is
not part of the five-rung ladder.

\textbf{Two elements of Definition 2.4 were formalised after the Section
3 scores had been computed:} the 5\% tie band, and the comparator
declarations for the non-nested rungs. This is the substantive
qualification on the protocol, stated in the main text. Neither element
changes any retention outcome: the tie band never decides a retention
result, the smallest deficit being 17\%.

\textbf{Specification B was run last and was never pooled with the NCAM
2016 domain.} The constraint is recorded in the code itself, whose
module docstring reads ``Wave E ladder on Ω\_xte alone. Do not pool with
NCAM 2016.'', and in \texttt{meta.json}, which records that the full
xteNCAM SSB table was not extracted and was not pooled. The two
specifications are therefore scored separately throughout, and no claim
in the main text transfers a result from one to the other.

\begin{center}\rule{0.5\linewidth}{0.5pt}\end{center}

\subsection{\texorpdfstring{SI-2 Schaefer is not the (\mathfrak s \to 0)
limit of the depensation
family}{SI-2 Schaefer is not the (s ) limit of the depensation family}}\label{si-2-schaefer-is-not-the-s-limit-of-the-depensation-family}

\textbf{Statement.} In the family of Definition 2.1, the Schaefer law
corresponds to the choice (a \equiv 1). The limit (\mathfrak s \to 0) in
the depensation branch yields (a(S\_t) = S\_t/K), hence a cubic
modification of the logistic surplus, and not the Schaefer law.

\textbf{Proof.} Substituting (\mathfrak s = 0) in (a(S\_t) = (S\_t -
\mathfrak s)/(K - \mathfrak s)) gives (a(S\_t) = S\_t/K), so (g(S\_t) =
r S\_t (1 - S\_t/K)\cdot(S\_t/K)), a cubic surplus term. The Schaefer
law requires (a \equiv 1), which is the separate branch in which the
depensation factor is replaced, not obtained by a limiting value of
(\mathfrak s). ∎

The distinction matters because it fixes what M1b adds to M1: a separate
branch of the family, not a limiting case of it. Abaee (2026b) makes the
same distinction.

\begin{center}\rule{0.5\linewidth}{0.5pt}\end{center}

\subsection{SI-3 Uncertainty layer: bandwidth grids, disagreement rows,
and
provenance}\label{si-3-uncertainty-layer-bandwidth-grids-disagreement-rows-and-provenance}

\subsubsection{Bandwidth and block-length
sensitivity}\label{bandwidth-and-block-length-sensitivity}

The M1-against-persistence margin was recomputed over a grid of HAC
truncation lags and moving-block bootstrap block lengths, because the
automatic (h-1) truncation and the (\max(h,3)) block length are not
neutral choices when origins overlap at (h=5), the estimation window
expands, the predictand is a smoothed reconstruction, and the sample
straddles a structural break.

{\def\LTcaptype{none} % do not increment counter
\begin{longtable}[]{@{}lll@{}}
\toprule\noalign{}
Quantity & Specification A & Specification B \\
\midrule\noalign{}
\endhead
\bottomrule\noalign{}
\endlastfoot
DM (z), (h=1), lag 0 → lag 3 & 1.14 → 2.30 & 2.81 → 1.70 \\
Bootstrap (p), (h=1), block 3--9 & 0.17--0.19 & 0.036--0.049 \\
Bootstrap (p), (h=5), block 3--9 & 0.26--0.28 & 0.008--0.026 \\
\end{longtable}
}

The DM statistic crosses the conventional 1.96 threshold within the
range of defensible bandwidths on both specifications --- upward on A,
downward on B --- so no verdict rests on it. The bootstrap is markedly
more stable. The qualitative reading in the main text (Specification A
within noise, Specification B separating) is therefore a bootstrap
result rather than a DM result. The grid was run for the
M1-against-persistence margin; robustness across block lengths is
established for that comparison and is not asserted for every module.

\subsubsection{The four disagreement
rows}\label{the-four-disagreement-rows}

Four of the twenty-eight rows of Table 9 have a bootstrap interval
excluding zero while (\textbar z\textbar{} \textless{} 1.96):

{\def\LTcaptype{none} % do not increment counter
\begin{longtable}[]{@{}llllll@{}}
\toprule\noalign{}
Spec & (h) & Module & Comparator & DM (z) & 95\% CI (kt) \\
\midrule\noalign{}
\endhead
\bottomrule\noalign{}
\endlastfoot
A & 1 & M4 & M3 & 0.99 & {[}+4.7, +144.7{]} \\
B & 1 & M3 & persist & 1.85 & {[}+1.0, +92.5{]} \\
B & 5 & M1b & persist & 1.80 & {[}+18.2, +250.9{]} \\
B & 5 & M4 & M3 & 1.88 & {[}+20.2, +177.4{]} \\
\end{longtable}
}

Both summaries are reported side by side and no verdict rests on either
where they disagree.

\subsubsection{Provenance}\label{provenance}

The layer is produced by \texttt{campaign\_e1\_dm\_uncertainty.py},
deterministic under seed 0, with output archived at
\texttt{results/e1\_dm\_uncertainty.csv}.

\begin{center}\rule{0.5\linewidth}{0.5pt}\end{center}

\subsection{SI-4 Log-floor binding
counts}\label{si-4-log-floor-binding-counts}

The log-RMSE scores of Table 4 use (\log\max(\hat S,
\varepsilon\emph{\{\mathrm{log}\})) with (\varepsilon}\{\mathrm{log}\} =
10\^{}\{-3\}) kt, the trajectory code clipping the state to
({[}\varepsilon\_\{\mathrm{log}\}, 10\^{}\{6\}{]}) kt. The floor binds
on the following counts of origins, taken from the archived per-origin
records (Specification A: the annual-landings rolling pass of Section
3.2).

{\def\LTcaptype{none} % do not increment counter
\begin{longtable}[]{@{}llllll@{}}
\toprule\noalign{}
Spec & (h) & origins & M1 & M1b & M2, M3, M4 \\
\midrule\noalign{}
\endhead
\bottomrule\noalign{}
\endlastfoot
A & 1 & 25 & 15 & 17 & M3 at 3; M2 and M4 at most once per horizon \\
A & 5 & 21 & 19 & 19 & M3 at 3; M2 and M4 at most once per horizon \\
B & 1 & 59 & 22 & 24 & between 0 and 11 \\
B & 5 & 55 & 36 & 46 & between 0 and 11 \\
\end{longtable}
}

The raw-RMSE column, not the log column, is the retention score, so
these counts qualify a reported diagnostic rather than the verdict.

\begin{center}\rule{0.5\linewidth}{0.5pt}\end{center}

\subsection{SI-5 Retention-rule operating characteristics: full
decomposition}\label{si-5-retention-rule-operating-characteristics-full-decomposition}

Design and results are archived in full as
\texttt{E1\_SIMULATION\_RESULTS.md}, with the pre-registration at
\texttt{wave\_e\_cod/SPECIFICATION\_v4.md} (locked before execution) and
the per-replicate output at
\texttt{wave\_e\_cod/results/sim\_retention\_power.csv} (10,000 rows).
This section records the detail condensed out of Section 3.7.

\subsubsection{SI-5.1 Where the power is
lost}\label{si-5.1-where-the-power-is-lost}

H2 alone is ``beats persistence by more than the 5\% band at both
horizons'', ignoring the comparator gate.

{\def\LTcaptype{none} % do not increment counter
\begin{longtable}[]{@{}
  >{\raggedright\arraybackslash}p{(\linewidth - 10\tabcolsep) * \real{0.1667}}
  >{\raggedright\arraybackslash}p{(\linewidth - 10\tabcolsep) * \real{0.1667}}
  >{\raggedright\arraybackslash}p{(\linewidth - 10\tabcolsep) * \real{0.1667}}
  >{\raggedright\arraybackslash}p{(\linewidth - 10\tabcolsep) * \real{0.1667}}
  >{\raggedright\arraybackslash}p{(\linewidth - 10\tabcolsep) * \real{0.1667}}
  >{\raggedright\arraybackslash}p{(\linewidth - 10\tabcolsep) * \real{0.1667}}@{}}
\toprule\noalign{}
\begin{minipage}[b]{\linewidth}\raggedright
DGP
\end{minipage} & \begin{minipage}[b]{\linewidth}\raggedright
truth
\end{minipage} & \begin{minipage}[b]{\linewidth}\raggedright
passes H2 alone
\end{minipage} & \begin{minipage}[b]{\linewidth}\raggedright
passes full rule
\end{minipage} & \begin{minipage}[b]{\linewidth}\raggedright
gates, absolute
\end{minipage} & \begin{minipage}[b]{\linewidth}\raggedright
gates, conditional on H2
\end{minipage} \\
\midrule\noalign{}
\endhead
\bottomrule\noalign{}
\endlastfoot
D1 & M1 & 0.975 & 0.975 & 0.000 & 0\% removed \\
D2 & M1 & 0.420 & 0.420 & 0.000 & 0\% \\
D3 & M2 & 0.328 & 0.100 & 0.228 & 69\% \\
D4 & M1b & 0.180 & 0.010 & 0.170 & 94\% \\
\end{longtable}
}

The absolute cost of the comparator gates is therefore 0.17--0.23 of the
shortfall on the two cells where they bind.

Both framings hold. In absolute terms the dominant failure is H2 --- the
true module often does not out-predict persistence on data it generated.
Conditional on clearing H2, the comparator requirement is a dominant
further filter.

\subsubsection{SI-5.2 The chance
baseline}\label{si-5.2-the-chance-baseline}

With five structural modules, random assignment would place the
generating module first 20\% of the time.

{\def\LTcaptype{none} % do not increment counter
\begin{longtable}[]{@{}lll@{}}
\toprule\noalign{}
DGP & true module has lowest h=1 RMSE & versus chance \\
\midrule\noalign{}
\endhead
\bottomrule\noalign{}
\endlastfoot
D1 (M1 true) & 0.627 (62.7\%) & far better \\
D2 (M1 true) & 0.645 (64.5\%) & far better \\
\textbf{D3 (M2 true)} & \textbf{0.0375} & \textbf{worse than chance} \\
D4 (M1b true) & 0.258 (25.8\%) & barely better \\
\end{longtable}
}

On D3 the stock-flow module's mean one-year error (85.7 kt) exceeds that
of the residual (77.9 kt) and Allee (79.9 kt) modules. Estimation noise
actively disadvantages the correct structure relative to its siblings at
this sample size.

\subsubsection{SI-5.3 Alternative decision
rules}\label{si-5.3-alternative-decision-rules}

Computed on the same archived replicates, no refitting.

{\def\LTcaptype{none} % do not increment counter
\begin{longtable}[]{@{}lll@{}}
\toprule\noalign{}
Decision rule & mean power & specificity \\
\midrule\noalign{}
\endhead
\bottomrule\noalign{}
\endlastfoot
Retention rule as adopted (H1--H3, 5\% band) & 0.376 & 0.978 \\
Without the comparator gate (H2, H3 only) & 0.476 & 0.972 \\
Persistence only, h=1, 5\% band & 0.562 & 0.955 \\
Persistence only, any margin & 0.542 & 0.765 \\
Full rule, no tie band & 0.436 & 0.772 \\
Full rule, 10\% band & 0.337 & 0.998 \\
MASE \textless{} 1 against the naive benchmark & 0.651 & 0.675 \\
Information criterion, \texttt{n·log\ MSE\ +\ 2k} & 0.509 & 0.992 \\
\end{longtable}
}

False-retention across the four structural DGPs is 0.044 per
module-replicate pair, or 0.176 falsely retained modules per replicate;
under the persistence-true null the per-module rate is 0.005.

\subsubsection{SI-5.4 Two failed pre-check
candidates}\label{si-5.4-two-failed-pre-check-candidates}

A diagnostic that told an analyst in advance whether the rule has power
on a given series would be more valuable than any rule variant. Two were
examined and neither works.

\begin{enumerate}
\def\labelenumi{\arabic{enumi}.}
\tightlist
\item
  \textbf{Dispersion of scores across modules.} Median coefficient of
  variation against true-module power: Spearman ρ = 0.52. Too weak to
  license a protocol.
\item
  \textbf{Margin of the best-scoring module over persistence.} Spearman
  ρ = 0.88, Pearson r = 0.94 --- strongly predictive, but it fails on
  two grounds. It is computed from the same out-of-sample scores the
  retention rule consumes, so it cannot gate the decision; and
  thresholding it at the tie band misclassifies both D3 cells, flagging
  ``apply'' where the generating module wins less often than chance.
\end{enumerate}

Identifying when the rule has power therefore remains an open problem.

\end{document}
